
\documentclass{article}

\usepackage{graphicx}
\usepackage{comment}
\usepackage[hidelinks]{hyperref}
\usepackage{tabularx}
\usepackage{booktabs}
\usepackage{makecell}
\usepackage{multirow}
\usepackage{amsmath}
\graphicspath{./Pictures/}

\newcolumntype{n}{>{\hsize=0.7\hsize \raggedright\arraybackslash}X}
\newcolumntype{w}{>{\hsize=1.6\hsize \raggedright\arraybackslash}X}
\newcolumntype{a}{>{\hsize=1.2\hsize \raggedright\arraybackslash}X}
\newcolumntype{b}{>{\hsize=0.8\hsize \raggedright\arraybackslash}X}

\makeatletter
\renewcommand\paragraph{\@startsection{paragraph}{4}{\z@}%
	{-2.5ex\@plus -1ex \@minus -.25ex}%
	{1.25ex \@plus .25ex}%
	{\normalfont\normalsize\bfseries}}
\makeatother

\begin{document}
	% Title Page
	\begin{titlepage}
   \begin{center}
       \vspace*{1cm}
        \Huge
       \textbf{Education plan}
        
        \Large
       \vspace{0.5cm}
        -How do we assure a broad knowledge base within our company
            
       \vspace{1.5cm}

       \textbf{Martin Friberg}
\begin{center}
\tiny
\begin{tabular}{ | m{5em} | m{5em}| m{10em} |m{5em}| m{5em} |m{5em} |  } 
\hline
Version Number& Published Date & Description of revision & Author & Approved by \\ 
\hline
1.0 & 2020-09-22 & Sections about Workshops, Presentations, Resource sharing, Group meetings and discussions, Communication directives & Martin Friberg & Patrik Palmgren \\
\hline
2.0 & 2020-11-18 & Education to be done, education history and rewrote all the other sections. How to follow up on education & Martin Friberg & N/A \\
\hline


\end{tabular}
\end{center}

     \vspace{1cm}
     
    \small November 2020
     


      \vfill
\includegraphics[width=\linewidth]{Pictures/logo_heartbyte_transparent_v_1_1 (1)}

    \vfill
            
   \end{center}
\end{titlepage}
	\clearpage
	
	% Change Log
	
{\Large Document Change Record}

\begin{table}[h]
	\noindent\makebox[\textwidth]{
		\begin{tabularx}{1.2\textwidth}{l|l|w|n|n}
			\toprule
			\makecell[l]{Version \\ Number} & 
			\makecell[l]{Published \\ Date} & Description of Revision & Author & Approved by \\
			\midrule
			1.0 & 2020/10/15 &  Initial version &Emma Johansson & \\
			1.1 & 2020/11/12  & Add SQA activities, Document Review and Inspection process	& Emma Johansson & \\
			\bottomrule
		\end{tabularx}
	}
\end{table}


	\clearpage
	
	% Table of Contents
	\tableofcontents
	\clearpage
	
	% Introduction: Why is this an important document?
	\section{Introduction}
	This Software Quality Assurance Plan (SQAP) describes the process, methods, standards and procedures that will be used to ensure the desired quality level of the HeartByte project.
	
	\subsection{Scope}
	This plan establishes the SQA activities performed throughout the life cycle of the project. The goal of the SQA process is to verify that all software and documentation to be delivered meet all technical requirements and the expectations of both the customer Region Östergötland (RÖ) and CEO Daniel Ståhl. The SQA procedures defined herein shall be used to examine all deliverable software and documentation to determine compliance with tehcnical and performance requirements. 
	

	\subsection{SQA Roles}
	Emma Johansson as Quality Coordinator is responsible of conducting the SQAP and its compliance within the project. The company's Document Manager Patrik Palmgren will assist in review and inspection processes. Testing methodologies are conducted by Gustav Karlsson, Test Leader. 
	
	All members of HeartByte are responsible to deliver work of good quality and identify shortcomings in both processes and work products.
	
	%\clearpage
	%\section{References}
	%This chapter includes a reference list to relevant documents. 
	%document number, name, issue, issue date
	
	
	\clearpage
	\section{Software Configuration Management}
	The purpose of Software Configuration Management (SCM) is to implement a controlled change process and define build and process management. The project uses GitLab as version control tool. The framework on how GitLab is used within the project is found in \textit{Framework for programming TDDC88}. The main principles specified to assure high quality code is the use of several levels of feature branches and merge requests. Each function are to be developed on a separate branch and requested to merge to a higher level feature branch when finished. The merge request is then to be audited by the responsible person for each level before being accepted and merged to the highest level of development feature branches. This process work to avoid defects being imported into the running code.
	
	To assure readability of code and traceability to features connected to requirements, processes for generating readme-files and commenting of code are specified (see \textit{Framework for programming TDDC88}).
	
	Continuous integration refers to the practice where developers integrate their work continuously. This methodology is used in HeartByte's project in order to improve predictability and easier troubleshooting. By continuously integrating changes there will be no unexpected grand errors when integrating modules.

	\clearpage
	\section{SQA Activities}
	The following section describes tasks that lies within the responsibility of SQA. The tasks are dynamic and not occurring only once.  
	
	% software quality audits performed focusing on different aspects, e.g.traceability, tool usage, tool configuration, configuration identification, baselines, inspections, change control, test coverage, requirement development, building, loading,, reporting of compiler warnings, etc.
	
	\subsection{Conduct Appraisal processes}
	The project will throughout its life cycle produce a number of deliverables and work products. The appraisal process of these are to be specified and overseen by SQA in collaboration with responsible team leader. These processes are defined in section~\ref{method} \nameref{method}
		
	\subsection{Evaluate Software Tools}
	SQA shall evaluate software tools, both existing and planned, used for software development and support. This is due to ensure that existing tools are adequate for the purpose they are intended and whether the capabilities of the tools are needed. Planned tools are assessed in the same way, but also by determining feasibility in adopting the tool with regard to available knowledge and resources for education within HeartByte. This evaluation is done in consultation with representatives from the management team as well as the team proposing the tool.
	
	\subsection{Evaluate Project Planning and Compliance}
	The Project Plan is a work product under supervision of SQA. The Plan is assessed by ensuring that relevant activities are specified and planned, as well as management processes. SQA shall also evaluate that the project conducts the activities stated in the Project Plan. Prior to release, each major version is to be inspected according to the process described in  section~\ref{method} \nameref{method}. 
	
	%\subsection{Evaluate System Requirements Analysis Process}
	%Requirements analysis establishes a common understanding of the customer’s requirements between that customer and the software project team.  An agreement with the customer on the requirements for the software project is established and maintained.  This agreement is known as allocating system requirements to software and hardware.  Section 4 lists the system requirements documents.
	
	
	\subsection{Evaluate Software Configuration Compliance}
	SQA is responsible for following the compliance of the software configuration. This includes to asses the use of GitLab in comparison to specified guidelines as well as the readability and traceability of the code. Furthermore the continuous integration process is to be evaluated in terms of how frequent work is integrated at each level and the size of the integrated work pieces.
	
	\subsection{Evaluate Document Review and Inspection Process Compliance}
	SQA is responsible of ensuring that document review and inspection processes are followed for relevant documents. This includes documentation of the processes as well. 
	
	%\subsection{Software quality }
	%\subsection{Evaluate Coverage Rate of Tests}
	% Quality Factors and metrics
	
	\begin{comment}
	\clearpage
	
	\section{Software Quality Assurance Records} \label{doc}
	The aim of the software quality assurance record (SQR) is to report activities performed as part of the quality assurance work defined in the SQAP. The SQR shall identify the activity performed, persons involved and documents or data assessed. The resulting observations, suggestions, deviations and/or actions shall also be included. 
	
	The SQR's produced shall be gathered in a software quality assurance summary. There, all SQA activities performed are listed with references to respective SQR. The summary shall include a conclusion regarding the adherence to SQAP that is continuously updated.
	\end{comment}
	
	\clearpage
	\section{Tools, Technologies and Methodologies} \label{method}
	
	\subsection{Document Review and Inspection}
	Documents produced in the project are internally reviewed and inspected to  ensure that they contain relevant, correct and unambiguous information presented in an easily accessible way.  
	
	\subsubsection{Change Log}
	All documents that are to be reviewed or inspected shall have a document change log. This facilitates in tracking changes and verifying that those are approved by someone other than the author. Version Number provides a means to refer to a specific version. Published Date is the date when the certain version is locked. Description of Revision is a brief description of what has been changed/added/removed and, if applicable, which chapter it concerns. Author is responsible for the changes made to the certain version. Approved by refers to who has conducted a review or inspection, thus only applicable for primary number increments (e.g. 2.0, 3.0). See table \ref{tab:changeLog}.
	
	\begin{table}[h]
		\noindent\makebox[\textwidth]{
			\begin{tabularx}{1.2\textwidth}{l|l|w|n|n}
				\toprule
				\makecell[l]{Version \\ Number} & 
				\makecell[l]{Published \\ Date} & Description of Revision & Author & Approved by \\
				\midrule
				1.0	& 2020-10-23 &	Initial Version	& Emma Johansson & Patrik Palmgren \\
				1.1 & 2020-11-07 &	Document Review Process added (ch 6) & Emma Johansson &  \\
				\bottomrule
			\end{tabularx}
		}
		\caption{Example of Change Log}
		\label{tab:changeLog}
	\end{table}
	
	To each revision of a document, an author needs to be chosen among those who have contributed to the changes. This person is responsible for the state of the document and conducting the review process if a new version is to be released. Minor changes to single chapters shall be issued with a decimal increment, for example from 1.1 to 1.2. If the author considers a series of minor changes in total constitute a large difference from the latest reviewed version, the document manager shall be contacted for review or inspection. After such a process, the document is issued as a primary version with a primary number increment, i.e from 1.7 to 2.0.  
	
	Each new version of a document shall be included in the change log. At the publishing date it should be added to Teams at applicable location in pdf format. This way each version of the documents are easy to find and locked for changes after approval.
	
	\subsubsection{Review Process}
	The review process shall be applied to all documents that are not strictly internal to a work group (e.g. management, development) or informal. Accordingly all documents intended for someone outside the work group, i.e. to be used in their work or as information, shall be reviewed. 
	That means that documents with notes or research not interesting for someone outside does not need to be reviewed. All other documents, like test results, guidelines and helping documents should be reviewed. When there are uncertainties, contact our Document Manager for discussion.
	
	When an author of a document wants to release the document as a new primary version, for example 2.0, the author needs to contact the Document Manager or Quality Coordinator (before the version is released by being published to Teams). The review is conducted by assigned reviewee who assess the document and return with eventual comment. If something needs to be changed, the author is responsible for this, and within three days the updated document should be handed back to the reviewee. When the document is approved, the reviewee signs the changelog and the document is converted to a pdf before being published on Teams at applicable location.
	
	\subsubsection{Inspection Process}
	The inspection process should be applied to all documents intended for external stakeholders of the project and especially the living documents. This includes all documents that are graded by the course management.
	The roles of the process are as follows:
	
	\smallskip
	\begin{table}[h]
		\begin{tabularx}{\textwidth}{lX}
			Inspection Leader & Responsible for planning and coordinating the process \\
			Recorder & Responsible of documenting defects, decisions and recommendations \\
			Reader & Leads the walk-through of the document at the inspection meeting \\
			Author & Responsible for performing rework to meet the inspection exit criteria \\
			Inspector & Identifies and describes defects. All attendants of Inspection Meeting are Inspectors \\
			Specialist Inspector & Someone with specific knowledge relevant for the document to inspect \\
		\end{tabularx}
	\end{table}
	
	
	The role of inspection leader and recorder is taken by Document Manager or Quality Coordinator. This person can be the reader as well, otherwise the inspection leader chooses someone else to fill the role. For each living document, a specialist inspector is chosen, see table \ref{tab:livedocsroles}. The inspection leader can invite more inspectors if needed. All attendants to an inspection meeting also performs the role as inspectors. These roles all need to be represented at the inspection meeting. The author may also be present.
	
		
	\smallskip
	\begin{table}[h]
		{\renewcommand{\arraystretch}{1.8}
		\begin{tabularx}{\textwidth}{lab}
			\toprule
			Document & \multicolumn{2}{l}{Inspector} \\
			\cline{2-3}
					 & Role	& Name \\
			\midrule
			Architecture Notebook & Lead Developer & Filip Eriksson \\
			\makecell[l]{Software Requirement \\ Specification} & Lead Tester & Gustav Karlsson \\
			Education Plan & Configuration Manager & Sam Anlér \\
			Project Plan & Strategic Product Manager & Max Klasson \\
				& \& Line Manager & Patrik Palmgren \\
			\makecell[l]{Software Quality \\ Assurance Plan} & Project Manager & Axel Trolme \\
			Test Plan &	Quality Coordinator	& Emma Johansson \\
			\bottomrule
		\end{tabularx}}
		\label{tab:livedocsroles}
	\end{table}
	
	The inspection process starts with the author of the document contacting the Document Manager (before the version is released by being published to Teams). An inspection leader for the current inspection is chosen, who contacts concerned parties to schedule the inspection. The inspection leader hands out the material, i.e. the document to be inspected and possibly related work products, and gives a brief introduction if needed. The inspectors, including the inspection leader, perform an individual checking of the document. All anomalies and defects, as well as ideas for improval or missing parts, shall be written down and handed to the inspection leader at least the day before the following inspection meeting. The inspectors shall be especially observant of things that are hard to understand, incorrect or unambiguous. The inspector needs to see the document from the perspective of one's role and evaluate if any relevant information is missing or need to be further specified. During the inspection meeting the reader presents the product and defects from all inspectors are  brought up chronologically. The defects shall be gathered in a single defect list. The inspection meeting shall result in an exit decision where the document is (1) accepted with no further verification, (2) accepted with rework verification performed by one member or (3) not accepted and sent to reinspection. The result is then communicated to the author who resolves eventual defects. The inspection leader shall be informed when the rework is made to verify this and close the inspection process. Finally, the change log is signed by the inspection leader and the document published as a new primary version. The recorder is responsible for publishing documentation of the process including involved members, defect list, exit decision and date of closure.
	
	\subsection{Software Evaluation}	
	In order to provide high quality code in terms of understandability, reliability and testability the following methodologies have been implemented in HeartByte's development process.
	
	\subsubsection{Programming Guidelines}
	The programming guidelines conducted for the project can be found in \textit{Framework for Programming}. This includes the process of continuous integration using GitLab, expectations on the developers and the main processes to be used. Furthermore it specifies which naming conventions that are to be used as well as how to document the code, which is mainly done through commenting the code. 
	
	\subsubsection{Code Review}
	To ensure that the produced code follows the programming guidelines, all code that is pushed to our main development branch on GitLab is merged using merge requests. Thus, all code is reviewed by a responsible developer before accepted to the main branch. The assessment includes verifying that the code passes the CI pipeline, that the code is understandable and follows our conventions. The responsibility of accepting merge requests lies on the Integrator Emil Strömberg, who may temporarily delegate the responsibility to another developer. From iteration four, all developers are given permission to accept other developers merge requests. This change serves to minimize the time from the request being made to the time of acceptance. By dividing the responsibility evenly between the integration is not dependant on only one person, which might delay it.
	
	\subsection{Traceability} \label{traceability}
	In order to provide traceability between the requirements, developed functionality and tests, a work pipeline is established using the functionality of GitLab Issue Boards. Each functional requirement defined in the SRS is added to the board as an issue with appropriate labels to categorize them. The issue is named with the requirement's ID and name, and includes the description and category. The priority level is indicated using labels. The board lists issues separately with the labels "requirements functional", "TestCase", "Ready to develop", "Being developed" and "Test requirement". The issues are initially on the list "requirements functional", upon an issue is moved to "TestCase" when a tester starts to develop a test case for testing the requirement. In this stage, the ID of input and output objects are set according to the naming convention aforementioned. When the test case it developed, the tester moves the issue to the list "Ready to develop". Issues lie here waiting for a developer to get assigned to it and start developing corresponding functionality in the product. Then the issue is moved to "Being developed" until finished and progressed to "Test requirement". At that time a tester will perform the test and depending on the result, the issue will either be labelled "Done" and closed, or remain open as the bug is reported by creating a new issue labelled "Bug". If the test fails entirely, the issue is moved back to "Being developed".
	
	\subsection{Software Testing}
	The product developed by HeartByte is evaluated through software testing to ensure that functional requirements are met as well as bugs are identified. The full specification regarding testing is to be found in \textit{Test Plan 3.0}. The following section will provide a brief overview of the methodology and how the quality of the tests are to be assessed.
	

	\begin{comment}
		Dynamic testing - What TEsting team does
		Static testing - What I do
		Detect overcomplecity in code, find security errors (help source code be more secured when being deployed), enforce best coding practices (specific to code language, helps other understand code), can create project specific rules
		
		technical dept - translates as the implied cost for additional rework that can occur if at an early stage an easy but not efficient solution is chosen. In the future the easy code may restrict scalability. Might be measured by static testing tool
		
		sonarqube open source static analysis tools. 
		can detect tricky issues: bugs, code smells, security vulnerability,acticate rules needed (quality profiles),
	\end{comment}
	\subsubsection{Methodology}
	The testing will include integration tests, regression tests, system tests and user tests. The back-end of the product is mainly based on OpenEHR, which obstruct extensive unit testing. The majority of tests are designed as end to end tests, where the functionality is tested from the user's point of view through simulating human interaction with the website. 
	
	The integration tests are developed straight from the functional requirements of the product. These are handled in GitLab boards as test cases as described above, section \ref{traceability}. These tests are either automated or manual, depending on how measurable the output of the requirement is.
	
	Regression testing is performed by putting some of the automated integration tests in the Continuous Integration pipeline on GitLab. To withdraw most value from regression testing, a larger portion of tests would need to be included in the CI than what will be in the project. The limitation is due to time constraints within the project.
	
	System tests are performed by manually testing the system from the user's point of view by testing all elements on the user interface that are not covered by the integration tests. 
	
	User tests are performed several times during the development to ensure that the product really does capture the needs of the intended end users. This is done both on the prototype and the deployed system. The methods used for the user tests include Thinking-out-Loud, Lostness and System Usability Scale (SUS).
	
	%How to measure non-functional requirements?
	
	\subsubsection{Tools}
	The software testing uses the frameworks Jest and WebdriverIO, which in turn uses Jasmine. To make the integration tests Selenium is used with the two webdrivers Chrome and Firefox. Bugtracking is done through GitLab using issues labeled Bugs and a specific board. The environment and tools used for software testing are described in detail in the \textit{Test Plan}. 
	
	\subsubsection{Reporting Test Results}
	The test results of integration tests that are not in the CI are primarily reported on GitLab using the board described in \ref{traceability}. The same information about requirements, corresponding test cases, type of test and test results (exept of bugs) can be seen in the Requirement Traceability Matrix (RTM). There the tests are also divided into iterations, on which testing metrics will be calculated.
	
	Regression tests in the CI are reported on GitLab in the pipeline when commits passes it. 
	
	%System tests in Excel??
	
	User tests shall be documented with an explicit description of the method of the tests, the results of it and adjustments to the requirements and system that were taken accordingly.
	
	\subsubsection{Testing Quality}
	The type of tests performed makes it difficult to use standardized static metrics to evaluate different types of code coverage. The following metrics are to be used.
	
	\paragraph{Interface Coverage} 
	The metric describes how much of the system from the user's point of view that the tests cover. To compute the metric the number of interactive objects, such as buttons and links, there is in the user interface. It is also needed to count how many of those that are covered by some kind of test (exclusive user test). It is computated by dividing the number of objects covered by tests with the total number of objects and multiplying it with 100 making it a percentage.
	
	\begin{equation}
		\text{Interface Coverage} = \frac{\text{Number of Objects Covered}}{\text{Total Number of Objects}} \times 100
	\end{equation}
		
	\paragraph{Requirement Coverage}
	This metric is another way to describe how complete the tests are. It provides a high level view of what requirements have test coverage. However, this will not show any result of the test which could indicate the quality of the test cases if they find any bugs. It is computed by counting the number of requirements that are covered by test cases, dividing this with the total number of requirements and multiplying it with 100 to make it a percentage.
	
	\begin{equation}
		\text{Requirement Coverage} = \frac{\text{Number of Requirements Covered}}{\text{Total Number of Requirements}} \times 100
	\end{equation}
	
	\paragraph{Passed Test Cases Percentage}
	The metric shows the relation between the number of passed tests compared to the total number of tests executed. This metric is applied to integration tests. It is computed by counting the number of passed tests, dividing this with the total number of executed test and multiplying it with 100 to make it a percentage.
	
	\begin{equation}
		\text{Passed Test Cases Percentage} = \frac{\text{Number of Passed Tests}}{\text{Total Number of Tests Executed}} \times 100
	\end{equation}
	
	\paragraph{Failed Test Cases Percentage}
	The metric shows the relation between the number of failed tests compared to the total number of tests executed. This metric is applied to integration tests. It is computed by counting the number of failed tests, dividing this with the total number of executed test and multiplying it with 100 to make it a percentage.
	
	\begin{equation}
		\text{Failed Test Cases Percentage} = \frac{\text{Number of Failed Tests}}{\text{Total Number of Tests Executed}} \times 100
	\end{equation}

	\paragraph{Fixed Bugs Percentage}
	The percentage of fixed bugs describes to what proportion the detected bugs or defects are actually fixed. If the defects found by tests are not fixed, testing is somewhat redundant. It is computed by counting the number of fixed bugs, dividing this with the total number of bugs reported and multiplying it with 100 to make it a percentage. 
	
	\begin{equation}
		\text{Fixed Bugs Percentage} = \frac{\text{Number of Fixed Bugs}}{\text{Total Number of Bugs Reported}} \times 100
	\end{equation}

	\paragraph{Bugs per Test Case Percentage}
	This metric shows what percentage of test cases identifies bugs. A high value indicates accurate test cases, but also a large portion of defects in the system. A low value on the other hand might indicate too simple test cases. It is computed by counting the number of test cases with reported bugs, dividing this with the total number of test cases and multiplying it with 100 to make it a percentage.
	
	\begin{equation}
		\text{Bugs per Test Case Percentage} = \frac{\text{Number of Test Cases with Reported Bugs}}{\text{Total Number of Test Cases}} \times 100
	\end{equation}
	
	\clearpage
	\section{Problem Reporting}
	To encourage HearByte's members to report problems regarding quality or compliance to processes of the process plan, the reporting shall be as easy as possible. Therefore everyone are encouraged to contact the Quality Coordinator preferably through Teams. Depending on the problems nature, the team leader of concerned team will likely be involved. If there is a large issue a representative of the management team will be involved to produce a plan for solving the problem.
	
	
\end{document}